\begin{abstract}
Les \q{ponts} d'O.~Caramello transforment une propriété d'un topos, lue sur deux présentations différentes de ce topos, en un théorème qui relie ces deux présentations. Dans un exposé aux Mines Paris (2024), elle a proposé de confier à une machine la combinaison d'une base de données de telles présentations avec une base de données de telles lectures, afin d'engendrer des théorèmes. Nous exécutons à la main une petite version de ce générateur et la décrivons comme une pipeline logicielle: deux registres en entrée (cinq ponts démontrés en Lean~4 dans la bibliothèque \lp{} de lean-pool; de petites catégories dont les préfaisceaux sont des structures familières), puis jointure, spécialisation, dépliage, réécriture, vérification (en Lean et par une preuve directe) et filtrage par la littérature. Nous ne rapportons que les résultats que nous n'avons pas trouvés publiés: (1) les représentations d'un carquois dans les ensembles vérifient le tiers exclu faible $\neg p\vee\neg\neg p$ exactement quand chaque sommet a au plus une flèche sortante (les automates sur deux lettres ne le vérifient donc pas); (2) deux carquois n'ont des catégories de représentations ensemblistes équivalentes que s'ils sont isomorphes, et les points essentiels sont les sommets; (3) les représentations du carquois cyclique à $n$ sommets ont exactement $2^n+1$ sous-topos, obtenus en inversant un ensemble quelconque de flèches ou en écrasant tout; (4) les points essentiels et les sous-topos des actions de $T_n$, $PT_n$, $I_n$, $O_n$, $M_n(\mathbb F_q)$ sont respectivement au nombre de $n$, $n+1$, $n+1$, $n$, $n+1$ et de $n+1$, $n+2$, $n+2$, $n+1$, $n+2$; (5) un transfert de la condition d'Ore aux coins $eMe$; (6) le tiers exclu faible pour une douzaine de topos combinatoires. Les résultats connus sont consignés à part, pour l'étalonnage, et une dernière section discute les limites de la démarche. Une version anglaise accompagne ce texte.
\end{abstract}
